\scanpage{40}\mspage{39}
% L1–3
\uline{NB}\qquad Let $(X_\alpha)$ be a \inserted{finite} family of tame subsets of $L$, \deleted{mutual\textsuperscript{ly}} mutual\textsuperscript{ly} tame i.e such that
% L4
\[
S=\bigcup_\alpha X_\alpha^{\bullet}
\]
% L5–6
is a totally ordered subset~:
\[
t_1<t_2<\ldots<t_n
\]
% L7–12
Then the Boolean algebra generated by the $X_\alpha$ \deleted{is \unsure{consti|enti}} has as \deleted{components} \deleted{\unsure{elements} of the spectrum (when one interprets them as subsets of $L$,} ``components'' the following tame sets \inserted{(and pseudo-segment pieces)} (arranged in increasing order for $<$)
% L13–14
\[
\scriptstyle
\underset{\substack{\vert\\\inserted{\text{if }t_1\notin\partial L}}}{[L_{<t_1}]}\unsure{,|a}\;
\{t_1\},\ ]t_1,t_2[,\ \{t_2\},\ ]t_{1},t_3[,\ \{t_3\},\ldots,\
\deleted{\unsure{\{t\}}}\ \{t_n\},\
\underset{\substack{\rule{3em}{0.35pt}\\\inserted{\text{if }t_n\notin\partial L}}}{[L_{>t_n}]}
\]
% L15–18
Thus, this Boolean algebra is known when one knows $S$, i.e. the points $t_1<\ldots<t_n$ (\unsure{or the|or} ``decomposition'' \inserted{or ``subdivision''} of the segment piece \unsure{$L$|$\Lambda$|$A$}).
